% Chapter 9 of The Replica Rosetta Stone: ridge regression and double descent.
\rstchapter{Ridge Regression}{Three derivations of double descent}{replicas\dict random matrices and the CGMT}{ch:ridge}
\stonefoot{proved}
\begin{refsection}[chapters/09-ridge.bib]

% local notation: the limiting risk
\newcommand{\ridgeR}{\mathcal{R}}

\begin{chaptersummary}
Fit $n$ noisy linear observations in dimension $d$ by ridge regression and let $n,d\to\infty$ with $n/d\to\alpha$. The excess risk tends to an explicit function of $\alpha$ and the ridge $\lambda$; without ridge it diverges at the interpolation threshold $\alpha=1$, and with the optimal ridge it decreases in $\alpha$. Physicists obtain this learning curve from a replica-symmetric saddle point at zero temperature. We prove it completely by random matrices, where the replicas' susceptibility is the trace of the resolvent and the cavity field is the leave-one-out prediction, and in outline by Gordon's comparison, whose auxiliary problem has the same stationarity equations as the replica saddle point. A simulation at $n=1000$ agrees with the curves within its intervals; everything about the isotropic Gaussian case is proved, and what remains open concerns real data and nearby models.
\end{chaptersummary}

\section{The question: the test error of ridge regression when \texorpdfstring{$n\approx d$}{n = d}}\label{sec:ridge:question}

Least squares with an $\ell_2$ penalty is one of the oldest regularized estimators in statistics.\sidenote{The name is from \textcite{HoerlKennard1970}. In scikit-learn's \texttt{Ridge} the parameter \texttt{alpha} is $n\lambda$ in our normalization, and the default intercept must be switched off to match our model.} With $d$ fixed and $n$ growing, the excess risk of least squares is about $\sigma^2d/n$. With $n$ and $d$ of the same order it is not monotone in $n/d$: it blows up at $n=d$, where the data are first fitted exactly, and comes down on the other side. \textcite{BelkinEtAl2019} named the shape \term{double descent}\index{double descent}; for the linear model, physicists had drawn it nearly thirty years earlier \parencite{KroghHertz1991,LoogEtAl2020}.

\paragraph{The model.} We use the conventions of Part~III. The inputs $x_1,\dots,x_n\in\R^d$, the rows of the design $X\in\R^{n\times d}$, are independent $\Normal(0,I_d)$ vectors, and $y_i=\ip{x_i}{\beta^*}+\sigma\varepsilon_i$ with a teacher of norm $\norm{\beta^*}=r$ and noise of mean $0$ and variance $1$. The \term{ridge estimator}\index{ridge regression} with penalty $\lambda>0$ is
\begin{equation}\label{eq:ridge:estimator}
\hat\beta_\lambda=\argmin_{\beta\in\R^d}\Bigl\{\frac1{2n}\norm{y-X\beta}^2+\frac\lambda2\norm\beta^2\Bigr\}=(X^\top X+n\lambda I)^{-1}X^\top y,
\end{equation}
and its limit as $\lambda\downarrow0$ is the \term{minimum-norm interpolator}\index{minimum-norm interpolator} $\hat\beta_0=X^+y$, which fits the data exactly when $d\ge n$. For a fresh input $x\sim\Normal(0,I_d)$ the excess risk is $R(\hat\beta)=\E_x(\ip{x}{\hat\beta}-\ip{x}{\beta^*})^2=\norm{\hat\beta-\beta^*}^2$, the test error is $R+\sigma^2$, and the training error is $T(\hat\beta)=\frac1n\norm{y-X\hat\beta}^2$. The proportional limit takes $n,d\to\infty$ with $n/d\to\alpha$, the load,\sidenote{This is Gardner's $\alpha$ and Donoho's $\delta$; \textcite{HastieEtAl2022} and \textcite{DobribanWager2018} write $\gamma=d/n=1/\alpha$.} and the \term{interpolation threshold}\index{interpolation threshold} is $\alpha=1$.

\begin{physicsbox}[replica symmetry at zero temperature]
Let $\chi>0$ solve $\lambda\chi+\dfrac{\chi}{1+\chi}=\dfrac1\alpha$ and put $\nu=\dfrac{\alpha\chi^2}{(1+\chi)^2}$. In the proportional limit the excess risk and the training error of ridge regression are
\[
R=\frac{r^2(\alpha\lambda\chi)^2+\sigma^2\nu}{1-\nu},\qquad T=\frac{R+\sigma^2}{(1+\chi)^2}.
\]
As $\lambda\to0$, $R\to\sigma^2/(\alpha-1)$ for $\alpha>1$ and $R\to r^2(1-\alpha)+\sigma^2\alpha/(1-\alpha)$ for $\alpha<1$: the risk of the minimum-norm interpolator diverges at $\alpha=1$. The ridge $\lambda^*=\sigma^2/(\alpha r^2)$ is optimal, its risk is $\sigma^2\chi$, and its curve has no peak.
\end{physicsbox}

\Cref{fig:ridge:curves} draws the prediction. Below $\alpha=1$ the interpolator fits the noise exactly, and its risk rises to the peak as $\alpha\to1$, so that there more data hurt \parencite{Nakkiran2019}; a small ridge caps the peak, and the optimal ridge removes it. How do physicists derive these curves, and what is proved? The second answer is short: all of it, in the isotropic Gaussian case, by two rigorous routes whose steps can be set against the physicist's one by one.

\section{The physics: replicas at zero temperature}\label{sec:ridge:physics}

A physicist reads an estimator defined by minimization as the ground state of a Gibbs measure, and computes its typical properties from the free entropy.\sidenote{Since $\beta$ denotes the regression coefficients, we write $b$ for the inverse temperature.}

\begin{heuristic}[Replicas at zero temperature]
\emph{Set up.} In physics units the weights are $w=\sqrt d\,\beta$ and the teacher is $w^*=\sqrt d\,\beta^*$, so that $\norm{w^*}^2/d=r^2$. The local fields of sample $\mu$ are $\zeta_\mu=\ip{x_\mu}{w}/\sqrt d$ and $\zeta^*_\mu=\ip{x_\mu}{w^*}/\sqrt d$, and $y_\mu=\zeta^*_\mu+\sigma\varepsilon_\mu$. Multiplying the ridge objective \eqref{eq:ridge:estimator} by $n$ gives the energy
\begin{equation}\label{eq:ridge:energy}
H(w)=\frac12\sum_{\mu=1}^n(y_\mu-\zeta_\mu)^2+\frac{\alpha\lambda}2\norm w^2,
\end{equation}
whose Gibbs measure $e^{-bH(w)}\dd w/Z$ concentrates on $\sqrt d\,\hat\beta_\lambda$ as $b\to\infty$. For a sample $w$ put $Q=\norm w^2/d$ and $m=\ip{w}{w^*}/d$, and for two independent samples, two \term{replicas}, $q=\ip{w^1}{w^2}/d$; at $b=\infty$ the replicas coincide, and $R=q-2m+r^2$.

\emph{Replicate.} The free entropy $\frac1d\E\log Z$ is computed from $\E Z^k$ at integer $k$ through $\E\log Z=\lim_{k\to0}\frac1k\log\E Z^k$. Given $k$ replicas, the fields $(\zeta^*,\zeta^1,\dots,\zeta^k)$ of one sample are centered Gaussian with covariances $Q_{ac}=\ip{w^a}{w^c}/d$, $m_a=\ip{w^a}{w^*}/d$ and $r^2$, and the samples are independent. With $e^{dS(Q,m)}$ the volume of the replica configurations with prescribed overlaps and $I(Q,m)=\E\exp\{-\frac b2\sum_a(\zeta^*+\sigma\varepsilon-\zeta^a)^2\}$,
\[
\E Z^k=\int\exp\Bigl\{d\Bigl[S(Q,m)-\frac{b\alpha\lambda}2\sum_aQ_{aa}+\alpha\log I(Q,m)\Bigr]\Bigr\}\dd Q\dd m ,
\]
which for large $d$ is dominated by its saddle point. The components of the replicas orthogonal to $w^*$ have Gram matrix $d(Q-mm^\top/r^2)$, and the volume of $k$ vectors in $\R^{d-1}$ with Gram matrix $C$ is proportional to $\det(C)^{(d-k-2)/2}$ (the Wishart density), so $S=\frac12\log\det(Q-mm^\top/r^2)$ up to constants and terms of order $1/d$.

\emph{Replica symmetry.} Assume that the saddle point does not distinguish the replicas: $Q_{aa}=Q$, $Q_{ac}=q$ for $a\ne c$, $m_a=m$. Then $Q-mm^\top/r^2$ has the eigenvalue $Q-q$ with multiplicity $k-1$ and the eigenvalue $Q-q+k(q-m^2/r^2)$, so that $\frac1kS=\frac12\log(Q-q)+\frac{q-m^2/r^2}{2(Q-q)}+O(k)$. For the energy write $\zeta^a=\omega+\sqrt{Q-q}\,z_a$, with $(\omega,\zeta^*)$ centered Gaussian, $\E\omega^2=q$, $\E\omega\zeta^*=m$, and $z_1,\dots,z_k$ independent standard Gaussians; this reproduces the covariances, and $\frac1k\log I=\E_{\omega,y}\log\E_ze^{-\frac b2(y-\omega-\sqrt{Q-q}\,z)^2}+O(k)$. The shared part $\omega$ of the fields is the \term{cavity field}\index{cavity field}.

\emph{Zero temperature.} As $b\to\infty$ the replicas collapse onto the ground state, with $Q-q=\chi/b$ for a finite \term{susceptibility}\index{susceptibility} $\chi$, which a probabilist would read as $b$ times the variance per coordinate of the Gibbs measure. The $z$-average is dominated by its saddle point (for the square loss it is a Gaussian integral, and the formula is exact up to a factor $(1+\chi)^{-1/2}$):
\begin{equation}\label{eq:ridge:prox}
\begin{aligned}
\E_ze^{-\frac b2(y-\omega-\sqrt{\chi/b}\,z)^2}&\asymp\exp\Bigl\{-b\min_\zeta\Bigl[\frac{(\zeta-\omega)^2}{2\chi}+\frac{(y-\zeta)^2}2\Bigr]\Bigr\}\\
&=\exp\Bigl\{-\frac{b\,(y-\omega)^2}{2(1+\chi)}\Bigr\},
\end{aligned}
\end{equation}
the minimum being attained at $\hat\zeta=(\omega+\chi y)/(1+\chi)$. Keeping the terms of order $b$ and using $\E(y-\omega)^2=r^2+\sigma^2-2m+q$, we get $\frac1{dk}\log\E Z^k\approx b\,\Psi(m,q,\chi)$ with the replica-symmetric potential
\[
\Psi(m,q,\chi)=\frac{q-m^2/r^2}{2\chi}-\frac{\alpha\lambda q}2-\frac{\alpha(r^2+\sigma^2-2m+q)}{2(1+\chi)}.
\]
The ground-state energy per coordinate, $\min H/d$, is $-\Psi$ at the stationary point, where $\partial_q\Psi=\partial_m\Psi=\partial_\chi\Psi=0$:
\begin{equation}\label{eq:ridge:saddle}
\frac1\chi=\alpha\lambda+\frac\alpha{1+\chi},\qquad m=\frac{\alpha r^2\chi}{1+\chi},\qquad\frac{q-m^2/r^2}{\chi^2}=\frac{\alpha(r^2+\sigma^2-2m+q)}{(1+\chi)^2}.
\end{equation}
The training residuals are $y-\hat\zeta=(y-\omega)/(1+\chi)$, so the training error is $(R+\sigma^2)/(1+\chi)^2$.
\end{heuristic}

\begin{computation}{Solving the replica-symmetric equations}\label{comp:ridge:saddle}
\step \emph{The susceptibility.} Multiplying the first equation of \eqref{eq:ridge:saddle} by $\chi/\alpha$ gives $\lambda\chi+\chi/(1+\chi)=1/\alpha$, that is, $\alpha\lambda\chi^2+(\alpha\lambda+\alpha-1)\chi-1=0$, with one positive root
\[
\chi=\frac{1-\alpha-\alpha\lambda+\sqrt{(1-\alpha-\alpha\lambda)^2+4\alpha\lambda}}{2\alpha\lambda}.
\]
\step \emph{The overlap with the teacher.} The same equation reads $1-\alpha\lambda\chi=\alpha\chi/(1+\chi)$, so the second equation gives $m=r^2(1-\alpha\lambda\chi)$.
\step \emph{The risk.} Since $q=R+2m-r^2$, the third equation reads $R-(r^2-m)^2/r^2=\nu(R+\sigma^2)$ with $\nu=\alpha\chi^2/(1+\chi)^2$, and $(r^2-m)^2/r^2=r^2(\alpha\lambda\chi)^2$; hence $R(1-\nu)=r^2(\alpha\lambda\chi)^2+\sigma^2\nu$, the formula of the physics box.
\step \emph{No ridge.} Let $\lambda\to0$. If $\alpha>1$, step~1 gives $\chi/(1+\chi)\le1/\alpha$, so $\chi\le1/(\alpha-1)$, $\lambda\chi\to0$, $\chi\to1/(\alpha-1)$, $\nu\to1/\alpha$, $R\to\sigma^2/(\alpha-1)$ and $T\to\sigma^2(1-1/\alpha)$. If $\alpha<1$, it gives $\lambda\chi>1/\alpha-1>0$, so $\chi\to\infty$, $\alpha\lambda\chi\to1-\alpha$, $\nu\to\alpha$, $R\to r^2(1-\alpha)+\sigma^2\alpha/(1-\alpha)$ and $T\to0$.
\step \emph{The optimal ridge.} At $\alpha\lambda=\sigma^2/r^2$ try $q=m$. Then $R=r^2-m=\sigma^2\chi$, and the third equation becomes $\sigma^2\chi(1-\sigma^2\chi/r^2)=\nu\sigma^2(1+\chi)$, that is, $1-\alpha\lambda\chi=\alpha\chi/(1+\chi)$, which is step~2. So at $\lambda^*=\sigma^2/(\alpha r^2)$ the risk is $\sigma^2\chi$, and the overlap of two replicas equals their overlap with the teacher.
\end{computation}

\paragraph{What was not controlled.} The continuation of $\E Z^k$ from integer $k$ to $k\to0$, which the integer moments do not determine (\cref{ch:freeentropy}); the exchange of $d\to\infty$ with $k\to0$ and $b\to\infty$; Laplace's method as $k\to0$, where the number $k(k-1)/2$ of off-diagonal overlaps becomes negative and the saddle point need not be a maximum; the replica-symmetric ansatz; and the concentration of the overlaps. Here the energy is a strictly convex quadratic, so the Gibbs measure is a Gaussian centered at $\sqrt d\,\hat\beta_\lambda$, nothing can break the symmetry between two draws from it, and a physicist would call replica symmetry exact. The rest of the chapter shows that the answer is right, although no proof continues anything in $k$.

\paragraph{The optimal ridge as a Bayes estimator.} For the prior $\beta^*\sim\Normal(0,\frac{r^2}dI_d)$ and noise $\Normal(0,\sigma^2)$, the posterior mean of $\beta^*$ is $(X^\top X+\frac{\sigma^2d}{r^2}I)^{-1}X^\top y$, the ridge estimator at $\lambda^*$. So step~5 is the Nishimori line of \cref{ch:nishimori}: the Nishimori identity (\cref{lem:freeentropy:nishimori}) gives $m=q$, and a posterior mean minimizes the mean squared error. In physics units the optimal weight decay is $\alpha\lambda^*=\sigma^2/r^2$, whatever $\alpha$.

\paragraph{How the curve was first found.} Not with replicas. \textcite{HertzKroghThorbergsson1989} treated the learning dynamics of the one-layer perceptron with random patterns, and \textcite{KroghHertz1991} wrote the gradient dynamics of the linear perceptron with weight decay in the eigenbasis of the input correlation matrix, found the optimal decay $\sigma^2/\sigma_u^2$ (our $\alpha\lambda^*$), noted that it does not depend on the spectrum of that matrix, and plotted the learning curve against $\alpha$: without decay it has the peak at $\alpha=1$, with the optimal decay none.\sidenote{Their Figure~1 has $r^2=1$ and $\sigma^2=0.2$; our formula gives $0.358$ at $\alpha=1$ and $0.148$ at $\alpha=2$ for the optimally decayed curve, as the figure shows.} \textcite{KroghHertz1992} solved the dynamics analytically. These are averages over the spectrum of a random matrix. \textcite{DunmurWallace1993} then computed the linear perceptron with weight decay and noisy data by the Gardner (replica) method at finite temperature.

\section{The stone}\label{sec:ridge:stone}

The proof column holds both rigorous routes, random matrices (\cref{sec:ridge:proof}) and comparison (\cref{sec:ridge:cgmt}); in both, the susceptibility is a concrete random quantity and the replica equations reappear unchanged.

\begin{stone}{replicas at zero temperature}{random matrices; Gordon's comparison}
\stonestep{The object}{A Gibbs measure $e^{-bH(w)}$ at $b\to\infty$, and the overlaps $Q,q,m$ of its samples.}{The estimator $\hat\beta_\lambda=(X^\top X+n\lambda I)^{-1}X^\top y$, linear in $y$; no temperature, no free entropy.}
\stonestep{Average over the data}{$\E Z^k$ at integer $k$: the fields of each sample are Gaussian with covariance $(Q,m)$.}{Average the noise and the teacher's direction exactly (\cref{lem:ridge:quad,lem:ridge:rotation}); traces of $G=(X^\top X/n+\lambda I)^{-1}$ remain. Comparison: $v^\top Xw\to\norm wg^\top v+\norm vh^\top w$ (\cref{thm:gordon:cgmt}).}
\stonestep{Replica symmetry, $b\to\infty$}{$Q_{ac}=q+\delta_{ac}\chi/b$: the measure shrinks onto the ground state, with width $\chi=b(Q-q)$.}{Strict convexity: one minimizer. The Gibbs measure is Gaussian, and $b(Q-q)=\tr(X^\top X+n\lambda I)^{-1}=:\chi_n$ exactly (\cref{ex:ridge:gibbs}).}
\stonestep{The cavity field}{$\zeta^a=\omega+\sqrt{\chi/b}\,z_a$, and \eqref{eq:ridge:prox} gives $\hat\zeta=(\omega+\chi y)/(1+\chi)$.}{Leave one out (\cref{lem:ridge:loo}): $x_i^\top\hat\beta_\lambda=(x_i^\top\hat\beta_{(i)}+\tau_iy_i)/(1+\tau_i)$, $\tau_i\to\chi$; the cavity field is the leave-one-out prediction.}
\stonestep{$\partial_q\Psi=0$: the susceptibility}{$\lambda\chi+\dfrac{\chi}{1+\chi}=\dfrac1\alpha$.}{Trace of $\lambda G+GX^\top X/n=I$: $\lambda\chi_n+\frac1n\sum_i\ell_i=\frac dn$, with leverages $\ell_i\to\chi/(1+\chi)$ (\cref{prop:ridge:chi}). Comparison: $\partial_\theta\cG=\partial_a\cG=0$ with $1+\chi=\sqrt{a^2+\sigma^2}/\theta$ (\cref{prop:ridge:identity}).}
\stonestep{$\partial_m\Psi=\partial_\chi\Psi=0$: the risk}{$m=\alpha r^2\chi/(1+\chi)$, $q-\frac{m^2}{r^2}=\frac{\alpha\chi^2(R+\sigma^2)}{(1+\chi)^2}$; so $R(1-\nu)=r^2(\alpha\lambda\chi)^2+\sigma^2\nu$.}{$R\to-\alpha r^2\lambda^2\chi'+\sigma^2(\chi+\lambda\chi')$ with $\chi'=-\alpha\chi^2/(1-\nu)$: the same formula. Comparison: the direction of $w$ gives $m$, its component along $h$ gives $q-m^2/r^2$.}
\stonestep{The limits}{$k\to0$ after $d\to\infty$, then $b\to\infty$; then $\lambda\to0$ in the formula.}{No counterpart for $k\to0$. The exchange of $\lambda\to0$ and $n\to\infty$ is \cref{thm:ridge:risk}(ii): monotone comparison and an inverse Wishart mean.}
\end{stone}

\paragraph{What the translation shows.} Every step except the continuation in $k$ has a rigorous counterpart, and three are identities rather than analogies: the susceptibility is the trace of the \term{resolvent}\index{resolvent}, in the limit $1/\alpha$ times the Marchenko--Pastur Stieltjes transform at $-\lambda$; the cavity field is the leave-one-out prediction; and the replica equations are the stationarity equations of Gordon's scalar problem. The routes then part: the random-matrix proof uses the closed form of the estimator but extends to non-Gaussian designs, while the comparison proof needs Gaussian designs but works for every convex loss and penalty (\cref{ch:lasso,ch:classification}). And the physicist's ridgeless curve sets $\lambda=0$ after $d\to\infty$, an exchange of limits that the proof must justify and that at $\alpha=1$ matters at every finite size (\cref{sec:ridge:fails}).

\section{The proof: random matrices}\label{sec:ridge:proof}

\begin{assumption}[Isotropic Gaussian ridge regression]\label{ass:ridge:model}
For each $n$ let $d=d_n$ with $n/d_n\to\alpha\in(0,\infty)$. The design $X\in\R^{n\times d}$ has independent $\Normal(0,1)$ entries; the noise $\varepsilon\in\R^n$ is independent of $X$ and has independent entries with mean $0$, variance $1$ and $\E\varepsilon_i^4\le c_4$ for a constant $c_4$; the teacher $\beta^*\in\R^d$ is deterministic with $\norm{\beta^*}=r$ for a fixed $r\ge0$; $\sigma\ge0$ is fixed; and $y=X\beta^*+\sigma\varepsilon$.
\end{assumption}

\begin{theorem}[The proportional-limit ridge risk]\label{thm:ridge:risk}
Assume \cref{ass:ridge:model}. For $\lambda>0$ let $\chi=\chi_\alpha(\lambda)$ be the unique positive solution of
\begin{equation}\label{eq:ridge:chi}
\lambda\chi+\frac{\chi}{1+\chi}=\frac1\alpha,
\end{equation}
let $\nu=\alpha\chi^2/(1+\chi)^2$, and define
\begin{equation}\label{eq:ridge:risk}
\ridgeR(\alpha,\lambda)=\frac{r^2(\alpha\lambda\chi)^2+\sigma^2\nu}{1-\nu},\qquad\cT(\alpha,\lambda)=\frac{\ridgeR(\alpha,\lambda)+\sigma^2}{(1+\chi)^2}.
\end{equation}
Then $0<\nu<1$, and as $n\to\infty$ the following hold in probability.
\begin{enumerate}
\item For every $\lambda>0$, $R(\hat\beta_\lambda)\to\ridgeR(\alpha,\lambda)$ and $T(\hat\beta_\lambda)\to\cT(\alpha,\lambda)$.
\item If $\alpha>1$, then $R(\hat\beta_0)\to\sigma^2/(\alpha-1)$; if $\alpha<1$, then $R(\hat\beta_0)\to r^2(1-\alpha)+\sigma^2\alpha/(1-\alpha)$; if $\alpha=1$, then $R(\hat\beta_0)\to\infty$ when $\sigma>0$ and $R(\hat\beta_0)\to0$ when $\sigma=0$. In each case the limit equals $\lim_{\lambda\downarrow0}\ridgeR(\alpha,\lambda)$.
\item If $r>0$ and $\sigma>0$, then $\lambda\mapsto\ridgeR(\alpha,\lambda)$ is strictly decreasing on $(0,\lambda^*]$ and strictly increasing on $[\lambda^*,\infty)$, where $\lambda^*=\sigma^2/(\alpha r^2)$; moreover $\ridgeR(\alpha,\lambda^*)=\sigma^2\chi_\alpha(\lambda^*)$, and this optimally tuned risk is a strictly decreasing function of $\alpha$. If $\sigma=0<r$, $\lambda\mapsto\ridgeR(\alpha,\lambda)$ is strictly increasing, and if $r=0<\sigma$ it is strictly decreasing.
\end{enumerate}
\end{theorem}

The theorem is the physics box. In other normalizations, part~(i) is due to \textcite{Dicker2016} at $\lambda^*$, to \textcite{DobribanWager2018}, and to \textcite[Corollary~5]{HastieEtAl2022}; part~(ii) for $\alpha\ne1$ is Theorem~1 of \textcite{HastieEtAl2022}; the optimal ridge of part~(iii) and its risk are in all three. The proof below is assembled from standard tools so that each step answers a step of the replica computation.

\begin{remark}[The Marchenko--Pastur form]\label{rem:ridge:mp}
By \cref{prop:ridge:chi} below, $\chi_\alpha(\lambda)=\lim\frac1n\tr(X^\top X/n+\lambda I)^{-1}$; since the spectral distribution of $X^\top X/n$ converges to the \term{Marchenko--Pastur law}\index{Marchenko--Pastur law} $\mathrm{MP}$ with ratio $1/\alpha$ (\cref{thm:toolkit:mp}), its Stieltjes transform $m$ (\cref{def:toolkit:stieltjes}) satisfies $m(-\lambda)=\alpha\chi_\alpha(\lambda)$, and \eqref{eq:ridge:risk} reads
\[
\ridgeR(\alpha,\lambda)=\int\frac{r^2\lambda^2+\sigma^2s/\alpha}{(s+\lambda)^2}\,\mathrm{MP}(\dd s)=r^2\lambda^2m'(-\lambda)+\frac{\sigma^2}\alpha\bigl(m(-\lambda)-\lambda m'(-\lambda)\bigr),
\]
Corollary~5 of \textcite{HastieEtAl2022} with their $\gamma=1/\alpha$. With the \term{effective ridge}\index{effective ridge} $\lambda_{\mathrm{eff}}=\lambda(1+\chi)$ of \cref{ch:kernels}, $\ridgeR=(r^2\lambda_{\mathrm{eff}}^2+\sigma^2/\alpha)/((1+\lambda_{\mathrm{eff}})^2-1/\alpha)$, the flat-spectrum kernel learning curve.
\end{remark}

\subsection{Four lemmas}

The first lemma is the rigorous cavity: what removing one sample does to the estimator.\sidenote{It rests on the Sherman--Morrison formula $(B+xx^\top)^{-1}=B^{-1}-B^{-1}xx^\top B^{-1}/(1+x^\top B^{-1}x)$, valid for invertible $B$ with $1+x^\top B^{-1}x\ne0$.}

\begin{lemma}[Leave one out]\label{lem:ridge:loo}
Let $\lambda>0$, $A=X^\top X+n\lambda I$, and for a sample $i$ let $A_{(i)}=A-x_ix_i^\top$, $\tau_i=x_i^\top A_{(i)}^{-1}x_i$, and $\hat\beta_{(i)}=A_{(i)}^{-1}(X^\top y-y_ix_i)$, the ridge estimator computed without sample $i$ and with the same penalty $n\lambda$. Then the \term{leverage}\index{leverage} of sample $i$ and its fitted value are
\[
\ell_i=x_i^\top A^{-1}x_i=\frac{\tau_i}{1+\tau_i},\qquad x_i^\top\hat\beta_\lambda=\frac{x_i^\top\hat\beta_{(i)}+\tau_iy_i}{1+\tau_i},
\]
and $0\le\tr A_{(i)}^{-1}-\tr A^{-1}\le1/(n\lambda)$.
\end{lemma}

\begin{proof}
The Sherman--Morrison formula gives $A^{-1}=A_{(i)}^{-1}-A_{(i)}^{-1}x_ix_i^\top A_{(i)}^{-1}/(1+\tau_i)$. Applied to $x_i$ it gives $A^{-1}x_i=A_{(i)}^{-1}x_i/(1+\tau_i)$, hence $\ell_i=\tau_i/(1+\tau_i)$. With $b_i=X^\top y-y_ix_i$ we have $\hat\beta_\lambda=A^{-1}(b_i+y_ix_i)$, so $x_i^\top\hat\beta_\lambda=(A^{-1}x_i)^\top b_i+\ell_iy_i=(x_i^\top\hat\beta_{(i)}+\tau_iy_i)/(1+\tau_i)$. Taking traces, $\tr A_{(i)}^{-1}-\tr A^{-1}=x_i^\top A_{(i)}^{-2}x_i/(1+\tau_i)$, which lies between $0$ and $\norm{A_{(i)}^{-1}}\tau_i/(1+\tau_i)\le1/(n\lambda)$, because $A_{(i)}\succeq n\lambda I$.
\end{proof}

\begin{lemma}[Quadratic forms]\label{lem:ridge:quad}
Let $M$ be a fixed symmetric $k\times k$ matrix.
\begin{enumerate}
\item If $u$ is uniform on the unit sphere of $\R^k$, then $\E u^\top Mu=\tr M/k$ and $\Var(u^\top Mu)\le2\norm M^2/(k+2)$.
\item If $\xi\in\R^k$ has independent entries with mean $0$, variance $1$ and $\E\xi_j^4\le c_4$, then $\E\xi^\top M\xi=\tr M$ and $\Var(\xi^\top M\xi)\le(1+c_4)\tr M^2$; if $\xi$ is standard Gaussian, $\Var(\xi^\top M\xi)=2\tr M^2$.
\end{enumerate}
\end{lemma}

\begin{proof}
The plan is to expand the second moment for part~2 and to reduce part~1 to the Gaussian case. In $\E(\xi^\top M\xi)^2=\sum M_{jl}M_{st}\E\xi_j\xi_l\xi_s\xi_t$ only the terms whose indices match in pairs survive, so $\Var(\xi^\top M\xi)=2\sum_{j\ne l}M_{jl}^2+\sum_jM_{jj}^2(\E\xi_j^4-1)$, which is at most $\max(2,c_4-1)\tr M^2\le(1+c_4)\tr M^2$, because $c_4\ge(\E\xi_j^2)^2=1$, and equals $2\tr M^2$ when $\E\xi_j^4=3$. For part~1, write $u=z/\norm z$ with $z\sim\Normal(0,I_k)$, so that $u$ is independent of $\norm z$ and $z^\top Mz=\norm z^2u^\top Mu$. Taking means with $\E\norm z^2=k$ gives $\E u^\top Mu=\tr M/k$, and taking second moments with $\E\norm z^4=k(k+2)$ and $\E(z^\top Mz)^2=(\tr M)^2+2\tr M^2$ gives
\begin{align*}
\Var(u^\top Mu)&=\frac{(\tr M)^2+2\tr M^2}{k(k+2)}-\frac{(\tr M)^2}{k^2}\\
&=\frac2{k+2}\Bigl(\frac{\tr M^2}k-\Bigl(\frac{\tr M}k\Bigr)^2\Bigr)\le\frac{2\norm M^2}{k+2}.\qedhere
\end{align*}
\end{proof}

\begin{lemma}[Rotation invariance]\label{lem:ridge:rotation}
For $\lambda\ge0$ write $R_\lambda(X,\varepsilon;\beta^*)$ and $T_\lambda(X,\varepsilon;\beta^*)$ for the excess risk and the training error of $\hat\beta_\lambda$ computed from $y=X\beta^*+\sigma\varepsilon$. For every orthogonal matrix $O$, $R_\lambda(XO,\varepsilon;\beta^*)=R_\lambda(X,\varepsilon;O\beta^*)$, and the same holds for $T_\lambda$. Consequently, under \cref{ass:ridge:model}, the law of $R_\lambda(X,\varepsilon;\beta^*)$ is that of $R_\lambda(X,\varepsilon;ru)$ with $u$ uniform on the unit sphere and independent of $(X,\varepsilon)$, and the same holds for $T_\lambda$.
\end{lemma}

\begin{proof}
The plan is to check the identities for a fixed rotation and then to average over a random one. The response $y=XO\beta^*+\sigma\varepsilon$ is that of the design $X$ and the teacher $O\beta^*$, and the estimator from $(XO,y)$ is $O^\top\tilde\beta$, with $\tilde\beta$ the estimator from $(X,y)$, because $((XO)^\top XO+n\lambda I)^{-1}(XO)^\top=O^\top(X^\top X+n\lambda I)^{-1}X^\top$ and $(XO)^+=O^\top X^+$. Hence $\norm{O^\top\tilde\beta-\beta^*}=\norm{\tilde\beta-O\beta^*}$ and $\norm{y-XOO^\top\tilde\beta}=\norm{y-X\tilde\beta}$. Since $XO$ has the law of $X$ and is independent of $\varepsilon$, the law of $R_\lambda(X,\varepsilon;O\beta^*)$ is that of $R_\lambda(X,\varepsilon;\beta^*)$ for every fixed $O$; taking $O$ uniform on the orthogonal group and independent of $(X,\varepsilon)$ gives the last claim.
\end{proof}

\begin{lemma}[An inverse Wishart mean]\label{lem:ridge:wishart}
Let $Z\in\R^{N\times k}$ have independent $\Normal(0,1)$ entries, with $N\ge k+2$. Then $Z^\top Z$ is invertible almost surely and $\E\tr(Z^\top Z)^{-1}=k/(N-k-1)$.
\end{lemma}

\begin{proof}
The columns of $Z$ are almost surely linearly independent, since a Gaussian vector falls in a fixed proper subspace with probability $0$. Let $z_1$ be the first column and $P$ the orthogonal projection onto the span of the others. The Schur complement formula gives $((Z^\top Z)^{-1})_{11}=1/(z_1^\top(I-P)z_1)$, the inverse squared distance from $z_1$ to the other columns: the random-matrix cavity. Given the other columns, $z_1$ is standard Gaussian and $I-P$ is a projection of rank $N-k+1$, so $z_1^\top(I-P)z_1$ has the chi-square law with $m=N-k+1\ge3$ degrees of freedom, and
\[
\E\frac1{\chi^2_m}=\int_0^\infty\frac1x\,\frac{x^{m/2-1}e^{-x/2}}{2^{m/2}\Gamma(m/2)}\dd x=\frac{2^{m/2-1}\Gamma(m/2-1)}{2^{m/2}\Gamma(m/2)}=\frac1{m-2}.
\]
So $\E((Z^\top Z)^{-1})_{11}=1/(N-k-1)$, and the columns are exchangeable.
\end{proof}

\subsection{The susceptibility}

The next proposition is the rigorous form of the first replica equation, and its proof is the leave-one-out version of the cavity argument.\sidenote{The sum of the leverages, $\tr X(X^\top X+n\lambda I)^{-1}X^\top$, is the effective number of parameters of ridge regression; \eqref{eq:ridge:sumrule} below says that it equals $d-n\lambda\chi_n$.}

\begin{proposition}[The susceptibility is the trace of the resolvent]\label{prop:ridge:chi}
Let $X$ be as in \cref{ass:ridge:model}, $S=X^\top X/n$, and for $\lambda>0$ let $G=(S+\lambda I)^{-1}$,
\[
\chi_n(\lambda)=\frac1n\tr G=\tr(X^\top X+n\lambda I)^{-1},\qquad\psi_n(\lambda)=\frac1n\tr G^2=-\chi_n'(\lambda).
\]
Then $\E(\chi_n(\lambda)-\chi_\alpha(\lambda))^2\to0$ and $\psi_n(\lambda)\to-\chi_\alpha'(\lambda)$ in probability, and
\begin{equation}\label{eq:ridge:chiprime}
\chi_\alpha'(\lambda)=-\frac{\alpha\chi^2}{1-\nu},\qquad 1-\alpha\lambda\chi=\frac{\alpha\chi}{1+\chi},\qquad 0<\nu<1.
\end{equation}
\end{proposition}

\begin{proof}
The plan is to show that $\chi_n$ concentrates, that its mean nearly solves the self-consistent equation \eqref{eq:ridge:chi} because of the leave-one-out identity, and that the equation has one positive solution; the derivative then follows by convexity. Recall $\norm G\le1/\lambda$.

\emph{Concentration.} The function $f(X)=\tr G=n\chi_n$ is smooth, and $\dd G=-G(\dd S)G$ with $\dd S=(\dd X^\top X+X^\top\dd X)/n$ gives $\partial f/\partial X_{ik}=-\frac2n(XG^2)_{ik}$. Hence $\norm{\nabla f}^2=\frac4n\tr(G^2SG^2)\le\frac{4d}{n\lambda^3}$, because each eigenvalue $s/(s+\lambda)^4$ of $G^2SG^2$ is at most $\lambda^{-3}$. The Gaussian Poincaré inequality (\cref{thm:toolkit:poincare}) gives $\Var f\le\E\norm{\nabla f}^2$, so $\Var\chi_n\le4d/(n^3\lambda^3)$.

\emph{Leave one out.} Taking the trace of $\lambda G+GS=I$, with $GS=\frac1n\sum_iGx_ix_i^\top$ and $G=nA^{-1}$, \cref{lem:ridge:loo} gives
\begin{equation}\label{eq:ridge:sumrule}
\lambda\chi_n+\frac1n\sum_{i=1}^n\ell_i=\frac dn,\qquad\ell_i=\frac1nx_i^\top Gx_i=\frac{\tau_i}{1+\tau_i},
\end{equation}
where $\tau_i=\frac1nx_i^\top G_{(i)}x_i$ with $G_{(i)}=nA_{(i)}^{-1}$ independent of $x_i$. Let $c_n=\E\chi_n$. Given $G_{(i)}$, \cref{lem:ridge:quad} (Gaussian case) gives $\E(x_i^\top G_{(i)}x_i-\tr G_{(i)})^2=2\E\tr G_{(i)}^2\le2d/\lambda^2$; \cref{lem:ridge:loo} gives $0\le\tr G_{(i)}-\tr G\le1/\lambda$; and $\E\abs{\chi_n-c_n}\le(\Var\chi_n)^{1/2}$. Together,
\[
\E\abs{\tau_i-c_n}\le\frac{\sqrt{2d}}{n\lambda}+\frac1{n\lambda}+\frac{2\sqrt d}{n^{3/2}\lambda^{3/2}}=:\epsilon_n\to0.
\]

\emph{The equation.} The map $t\mapsto t/(1+t)$ is $1$-Lipschitz on $[0,\infty)$, so taking means in \eqref{eq:ridge:sumrule} gives $\abs{\varphi(c_n)-d/n}\le\epsilon_n$ for $\varphi(c)=\lambda c+c/(1+c)$. The function $\varphi$ is continuous and strictly increasing from $\varphi(0)=0$ to $\infty$, so it has a continuous inverse on $[0,\infty)$, the equation $\varphi(c)=1/\alpha$ has exactly one solution $\chi_\alpha(\lambda)>0$, and $c_n\to\chi_\alpha(\lambda)$ because $d/n\to1/\alpha$. With the concentration bound, $\E(\chi_n-\chi_\alpha)^2=\Var\chi_n+(c_n-\chi_\alpha)^2\to0$.

\emph{The derivative.} The second identity in \eqref{eq:ridge:chiprime} is \eqref{eq:ridge:chi} multiplied by $\alpha$; it gives $\alpha\chi/(1+\chi)<1$, hence $\nu=\frac{\alpha\chi}{1+\chi}\cdot\frac\chi{1+\chi}<1$. Differentiating \eqref{eq:ridge:chi} (the implicit function theorem applies, since $\varphi'>0$) gives $\chi'(\lambda+(1+\chi)^{-2})=-\chi$, and by \eqref{eq:ridge:chi}, $\lambda+(1+\chi)^{-2}=\frac1{\alpha\chi}-\frac\chi{(1+\chi)^2}=\frac{1-\nu}{\alpha\chi}$; this is the first identity. Finally, $\chi_n(\lambda)=\frac1n\sum_j(s_j+\lambda)^{-1}$, with $s_j$ the eigenvalues of $S$, is convex in $\lambda$, so for $0<h<\lambda$
\[
\frac{\chi_n(\lambda)-\chi_n(\lambda-h)}h\le\chi_n'(\lambda)\le\frac{\chi_n(\lambda+h)-\chi_n(\lambda)}h.
\]
Both bounds converge in probability to the difference quotients of $\chi_\alpha$, which tend to $\chi_\alpha'(\lambda)$ as $h\downarrow0$; hence $\chi_n'(\lambda)\to\chi_\alpha'(\lambda)$ in probability.
\end{proof}

\subsection{Proof of the theorem}

\begin{proof}[Proof of \cref{thm:ridge:risk}, part (i)]
The plan is to write the risk and the training error as quadratic forms in the noise and in the teacher, average them with \cref{lem:ridge:quad}, and recognize the traces that remain. By \cref{lem:ridge:rotation} we may assume that $\beta^*=ru$ with $u$ uniform on the unit sphere and independent of $(X,\varepsilon)$. With $S$, $G$, $\chi_n$ and $\psi_n$ as in \cref{prop:ridge:chi}, $\frac1n\tr(SG^2)=\frac1n\tr(G-\lambda G^2)=\chi_n-\lambda\psi_n$.

\emph{The risk.} Since $\hat\beta_\lambda=GX^\top y/n$ and $GS=I-\lambda G$, $\hat\beta_\lambda-\beta^*=-\lambda G\beta^*+\frac\sigma nGX^\top\varepsilon$, so $R(\hat\beta_\lambda)=B+C+V$ with $B=\lambda^2\beta^{*\top}G^2\beta^*$, $C=-\frac{2\lambda\sigma}n\beta^{*\top}G^2X^\top\varepsilon$ and $V=\frac{\sigma^2}{n^2}\varepsilon^\top XG^2X^\top\varepsilon$. Given $X$: part~1 of \cref{lem:ridge:quad} with $M=G^2$ gives $\E[B\mid X]=\lambda^2r^2\frac nd\psi_n$ and $\Var(B\mid X)\le2r^4/(d+2)$; the cross term has $\E[C\mid X,u]=0$ and $\E[C^2\mid X,u]=\frac{4\lambda^2\sigma^2}n\beta^{*\top}G^2SG^2\beta^*\le\frac{4\sigma^2r^2}{n\lambda}$; and part~2 with $M=XG^2X^\top/n^2$ gives $\E[V\mid X]=\sigma^2(\chi_n-\lambda\psi_n)$ and $\Var(V\mid X)\le(1+c_4)\sigma^4\frac1{n^2}\tr(S^2G^4)\le(1+c_4)\frac{\sigma^4d}{n^2\lambda^2}$, because $s^2/(s+\lambda)^4\le\lambda^{-2}$. By Chebyshev's inequality given $X$ and by \cref{prop:ridge:chi},
\begin{equation}\label{eq:ridge:riskderiv}
R(\hat\beta_\lambda)\to-\alpha r^2\lambda^2\chi'+\sigma^2(\chi+\lambda\chi')
\end{equation}
in probability. By \eqref{eq:ridge:chiprime}, $-\alpha r^2\lambda^2\chi'=r^2(\alpha\lambda\chi)^2/(1-\nu)$ and $\chi+\lambda\chi'=\chi(1-\nu-\alpha\lambda\chi)/(1-\nu)=\chi(\frac{\alpha\chi}{1+\chi}-\nu)/(1-\nu)=\nu/(1-\nu)$, so the limit is $\ridgeR(\alpha,\lambda)$.

\emph{The training error.} Let $K=XX^\top/n$. The push-through identity $X(X^\top X+n\lambda I)^{-1}=(XX^\top+n\lambda I)^{-1}X$ gives $y-X\hat\beta_\lambda=\lambda(K+\lambda I)^{-1}y$, so that $T(\hat\beta_\lambda)=\frac{\lambda^2}n\norm{(K+\lambda)^{-1}(X\beta^*+\sigma\varepsilon)}^2$. By the singular value decomposition $X^\top(K+\lambda)^{-2}X=nSG^2$, and $\frac1n\tr(K+\lambda)^{-2}=\psi_n+(1-\frac dn)\lambda^{-2}$ because $K$ and $S$ have the same nonzero eigenvalues. The same averaging as for the risk (bias mean $\lambda^2r^2\frac nd(\chi_n-\lambda\psi_n)$, cross term of mean $0$, conditional variances $O(1/n)$) gives
\[
T(\hat\beta_\lambda)\to\alpha r^2\lambda^2(\chi+\lambda\chi')+\sigma^2\Bigl(1-\frac1\alpha-\lambda^2\chi'\Bigr).
\]
Put $s=1/(1+\chi)$; by \eqref{eq:ridge:chiprime}, $\alpha\lambda\chi=1-\alpha(1-s)$ and $\nu=\alpha(1-s)^2$. The bias coefficient is $\alpha\lambda^2\nu/(1-\nu)=(\alpha\lambda\chi)^2s^2/(1-\nu)$, and expanding the squares, the noise coefficient is
\[
1-\frac1\alpha+\frac{\alpha\lambda^2\chi^2}{1-\nu}=\frac{(1-\frac1\alpha)(1-\alpha(1-s)^2)+\frac1\alpha(1-\alpha(1-s))^2}{1-\nu}=\frac{s^2}{1-\nu}.
\]
So $T(\hat\beta_\lambda)\to s^2(r^2(\alpha\lambda\chi)^2+\sigma^2)/(1-\nu)=s^2(\ridgeR+\sigma^2)=\cT(\alpha,\lambda)$.
\end{proof}

\begin{proof}[Proof of \cref{thm:ridge:risk}, part (ii)]
The plan is to split the risk of $\hat\beta_0$ into a projection of $\beta^*$ and a noise term that dominates the noise term of ridge, and to match the noise term's mean with \cref{lem:ridge:wishart}. Let $P=X^+X$, the orthogonal projection onto the row space of $X$, of rank $\min(n,d)$ almost surely. Since $\hat\beta_0=P\beta^*+\sigma X^+\varepsilon$ and $X^+\varepsilon$ lies in the row space, $R(\hat\beta_0)=\norm{(I-P)\beta^*}^2+\sigma^2\norm{X^+\varepsilon}^2$.

\emph{The bias.} If $n\ge d$ then $P=I$. If $d>n$, \cref{lem:ridge:rotation} with $\sigma=0$ shows that the bias has the law of $r^2\norm{(I-P)u}^2$ with $u$ uniform and independent of $X$, and given $X$, \cref{lem:ridge:quad} with $M=P$ gives mean $r^2(1-n/d)$ and variance at most $2r^4/(d+2)$. So the bias tends to $r^2(1-\alpha)_+$ in probability.

\emph{The noise.} If $n\ge d+2$, then $\E[\norm{X^+\varepsilon}^2\mid X]=\tr(X^\top X)^{-1}$, of mean $d/(n-d-1)$ by \cref{lem:ridge:wishart}; if $d\ge n+2$, then $X^+=X^\top(XX^\top)^{-1}$ and the mean is $\E\tr(XX^\top)^{-1}=n/(d-n-1)$. These tend to $N_0=1/(\alpha-1)$ if $\alpha>1$ and $N_0=\alpha/(1-\alpha)$ if $\alpha<1$. For a lower bound, write $X=\sum_js_ju_jv_j^\top$ over its positive singular values; then $X^+\varepsilon=\sum_js_j^{-1}(u_j^\top\varepsilon)v_j$, while $(X^\top X+n\lambda I)^{-1}X^\top\varepsilon=\sum_js_j(s_j^2+n\lambda)^{-1}(u_j^\top\varepsilon)v_j$ with $s_j/(s_j^2+n\lambda)<1/s_j$, so for every $\lambda>0$
\[
\norm{X^+\varepsilon}^2\ge Q_\lambda:=\norm{(X^\top X+n\lambda I)^{-1}X^\top\varepsilon}^2.
\]
The variance term of part~(i), with $\sigma=1$, shows $Q_\lambda\to N(\lambda):=\chi+\lambda\chi'=\nu/(1-\nu)$ in probability. As $\lambda\downarrow0$, step~4 of \cref{comp:ridge:saddle}, which uses only \eqref{eq:ridge:chi}, gives $\nu\to1/\alpha$ if $\alpha>1$ and $\nu\to\alpha$ if $\alpha<1$, so $N(\lambda)\to N_0$; if $\alpha=1$, \eqref{eq:ridge:chi} reads $\lambda\chi(1+\chi)=1$, so $\chi\to\infty$, $\nu\to1$ and $N(\lambda)\to\infty$.

\emph{Conclusion.} Let $\alpha\ne1$, $Y_n=\norm{X^+\varepsilon}^2$ and $\delta>0$, and choose $\lambda$ with $N(\lambda)>N_0-\delta/2$. Then $\Prob(Y_n<N_0-\delta)\le\Prob(Q_\lambda<N(\lambda)-\delta/2)\to0$. Since $\abs{Y_n-N_0}=(Y_n-N_0)+2(N_0-Y_n)_+$, $\E Y_n\to N_0$ and $\E(N_0-Y_n)_+\le\delta+N_0\Prob(Y_n<N_0-\delta)$, this lower bound and the matching mean give $\E\abs{Y_n-N_0}\to0$; with the bias, this is the claim for $\alpha\ne1$. If $\alpha=1$ and $\sigma>0$, then $R(\hat\beta_0)\ge\sigma^2Q_\lambda$ for every $\lambda$ and $N(\lambda)\to\infty$, so $R(\hat\beta_0)\to\infty$; if $\alpha=1$ and $\sigma=0$, $R(\hat\beta_0)$ is the bias, of mean $r^2(1-n/d)_+\to0$. Finally, $\ridgeR(\alpha,\lambda)=r^2(\alpha\lambda\chi)^2/(1-\nu)+\sigma^2N(\lambda)$, and the same limits show that its first term tends to $0$ if $\alpha>1$ and to $r^2(1-\alpha)$ if $\alpha<1$; if $\alpha=1$ it equals $r^2/(1+2\chi)\to0$, because there $\lambda\chi=1/(1+\chi)$ and $1-\nu=(1+2\chi)/(1+\chi)^2$.
\end{proof}

\begin{proof}[Proof of \cref{thm:ridge:risk}, part (iii)]
The plan is to differentiate \eqref{eq:ridge:riskderiv} and to show that one factor of the derivative has a fixed sign. Differentiating,
\[
\frac{\dd\ridgeR}{\dd\lambda}=-\alpha r^2\lambda(2\chi'+\lambda\chi'')+\sigma^2(2\chi'+\lambda\chi'')=(\sigma^2-\alpha r^2\lambda)(2\chi'+\lambda\chi'').
\]
To find the sign of $2\chi'+\lambda\chi''$, invert \eqref{eq:ridge:chi}: $\lambda=u-v$ with $u=1/(\alpha\chi)$ and $v=1/(1+\chi)$, a smooth function of $\chi$ with
\[
\lambda_\chi=-\frac u\chi+v^2=-\frac{1-\nu}{\alpha\chi^2}<0,\qquad\lambda_{\chi\chi}=\frac{2u}{\chi^2}-2v^3 .
\]
Then $\chi'=1/\lambda_\chi$, $\chi''=-\lambda_{\chi\chi}/\lambda_\chi^3$, and $2\chi'+\lambda\chi''=(2\lambda_\chi^2-\lambda\lambda_{\chi\chi})/\lambda_\chi^3$. Expanding,
\[
2\lambda_\chi^2-\lambda\lambda_{\chi\chi}=2uv\Bigl(\frac1\chi-v\Bigr)^2=\frac2{\alpha\chi^3(1+\chi)^3}>0,
\]
so $2\chi'+\lambda\chi''<0$ for every $\lambda>0$, and the derivative of $\ridgeR$ has the sign of $\alpha r^2\lambda-\sigma^2$, which proves the monotonicity in $\lambda$ in each case. At $\lambda^*$, $\alpha r^2\lambda^*=\sigma^2$, and \eqref{eq:ridge:riskderiv} gives $\ridgeR(\alpha,\lambda^*)=\sigma^2\chi+\lambda^*\chi'(\sigma^2-\alpha r^2\lambda^*)=\sigma^2\chi$. For the monotonicity in $\alpha$, put $c=\sigma^2/r^2$: at $\lambda=c/\alpha$, \eqref{eq:ridge:chi} multiplied by $\alpha$ reads $c\chi+\alpha\chi/(1+\chi)=1$, so $c\chi<1$ and $\alpha=(1-c\chi)(1+1/\chi)$, a product of two positive, strictly decreasing functions of $\chi\in(0,1/c)$ that maps $(0,1/c)$ onto $(0,\infty)$. Its inverse $\alpha\mapsto\chi_\alpha(\lambda^*)$ is therefore strictly decreasing, and so is $\sigma^2\chi_\alpha(\lambda^*)$.
\end{proof}

The \term{optimal ridge}\index{optimal ridge} is optimal at every finite size, which is Proposition~1 of \textcite{Dicker2016}: with $\beta^*=ru$, $\E[R(\hat\beta_\lambda)\mid X]=\frac1d\sum_j(r^2\lambda^2+\sigma^2\frac dns_j)/(s_j+\lambda)^2$ is minimized at $\lambda=\sigma^2d/(nr^2)$ for every design, with minimum $\sigma^2\chi_n$ (\cref{ex:ridge:dicker}).

\section{The proof: Gordon's auxiliary problem}\label{sec:ridge:cgmt}

The second rigorous route writes the ridge problem as a Gaussian min--max and replaces the design by two Gaussian vectors, with the \term{convex Gaussian min--max theorem}\index{convex Gaussian min-max theorem} (CGMT) of \cref{ch:gordon}.\sidenote{The idea of combining Gordon's comparison with convex duality to obtain exact asymptotics is Stojnic's (\cite{Stojnic2013d} for the LASSO); \textcite{ThrampoulidisOymakHassibi2015} and \textcite{ThrampoulidisAbbasiHassibi2018} turned it into the CGMT.} We compute the auxiliary problem as a physicist would, prove that its stationarity equations are the replica equations, and then say what the CGMT adds.

\begin{computation}{Gordon's auxiliary problem for ridge}\label{comp:ridge:ao}
\step \emph{The primary problem.} With $w=\beta-\beta^*$ the residual is $\sigma\varepsilon-Xw$, and $\frac1{2n}\norm{Xw-\sigma\varepsilon}^2=\max_{v\in\R^n}\{v^\top(Xw-\sigma\varepsilon)-\frac n2\norm v^2\}$, attained at $v=(Xw-\sigma\varepsilon)/n$. So the minimum of the ridge objective is
\[
\min_{w\in\R^d}\max_{v\in\R^n}\Bigl\{v^\top Xw-\sigma v^\top\varepsilon-\frac n2\norm v^2+\frac\lambda2\norm{\beta^*+w}^2\Bigr\},
\]
a bilinear form in the Gaussian matrix $X$ plus a function convex in $w$ and concave in $v$.
\step \emph{The auxiliary problem.} The CGMT compares it with the problem in which $v^\top Xw$ is replaced by $\norm wg^\top v+\norm vh^\top w$, with $g\sim\Normal(0,I_n)$ and $h\sim\Normal(0,I_d)$ independent of everything else.
\step \emph{The directions.} For $\norm v=\theta/\sqrt n$, the terms linear in $v$ are largest for $v$ parallel to $\norm wg-\sigma\varepsilon$, with value $\frac\theta{\sqrt n}\norm{\norm wg-\sigma\varepsilon}\approx\theta\sqrt{\norm w^2+\sigma^2}$. For $\norm w=a$, the terms linear in $w$, $w^\top(\theta h/\sqrt n+\lambda\beta^*)$, are smallest for $w$ antiparallel to $\theta h/\sqrt n+\lambda\beta^*$, with value $\approx-a\sqrt{\theta^2/\alpha+\lambda^2r^2}$, since $\norm h^2/n\approx1/\alpha$ and $\ip{h}{\beta^*}/\sqrt n\approx0$.
\step \emph{The scalar problem.} Exchanging the minimum over the direction of $w$ with the maximum over $\theta$, which the convexity of the limit allows, and adding $-\theta^2/2$ and $\frac\lambda2(r^2+a^2)$, the auxiliary problem reduces to $\min_{a\ge0}\max_{\theta\ge0}\cG(a,\theta)$ with
\begin{equation}\label{eq:ridge:ao}
\cG(a,\theta)=\theta\sqrt{a^2+\sigma^2}-a\sqrt{\theta^2/\alpha+\lambda^2r^2}-\frac{\theta^2}2+\frac\lambda2(r^2+a^2).
\end{equation}
\step \emph{Reading the solution.} At the saddle point $(a_*,\theta_*)$: $a_*=\norm{\hat\beta_\lambda-\beta^*}$, so $R=a_*^2$; $\theta_*^2=n\norm{v_*}^2=T$; the optimal direction gives $\ip{\hat\beta_\lambda}{\beta^*}=r^2-a_*\lambda r^2/\sqrt{\theta_*^2/\alpha+\lambda^2r^2}$; and the component of $\hat\beta_\lambda-\beta^*$ along $h$ has squared length $a_*^2(\theta_*^2/\alpha)/(\theta_*^2/\alpha+\lambda^2r^2)$.
\end{computation}

The vector $h$ plays the part of the Gaussian fluctuation of the replicas: the estimator is a multiple of $\beta^*$ plus an independent Gaussian vector, as the overlaps $m$ and $q-m^2/r^2$ describe. The correspondence is exact.

\begin{proposition}[Three systems, one solution]\label{prop:ridge:identity}
Let $\alpha,\lambda,r,\sigma>0$, let $\chi=\chi_\alpha(\lambda)$ and $\nu$ be as in \cref{thm:ridge:risk}, and let $\cG$ be the function \eqref{eq:ridge:ao} on $[0,\infty)^2$.
\begin{enumerate}
\item The function $\cG$ is strictly convex in $a$ and strictly concave in $\theta$, and it has exactly one stationary point $(a_*,\theta_*)$ in $(0,\infty)^2$; it is a saddle point, and
\[
1+\chi=\frac{\sqrt{a_*^2+\sigma^2}}{\theta_*},\qquad a_*^2=\ridgeR(\alpha,\lambda),\qquad\theta_*^2=\cT(\alpha,\lambda).
\]
\item The replica-symmetric equations \eqref{eq:ridge:saddle} have exactly one solution $(m,q,\chi)$ with $\chi>0$. Its $\chi$ is $\chi_\alpha(\lambda)$, $q-2m+r^2=\ridgeR(\alpha,\lambda)$, and $m$ and $q-m^2/r^2$ are the quantities read off the auxiliary problem in step~5 of \cref{comp:ridge:ao}.
\end{enumerate}
In particular the replica prediction, the auxiliary problem and \cref{thm:ridge:risk} give the same susceptibility, risk and training error.
\end{proposition}

\begin{proof}
The plan is to solve the stationarity equations of $\cG$ by introducing $\chi$ through the first formula of part~1. Convexity holds term by term: $\theta\sqrt{a^2+\sigma^2}$ is convex in $a$ and linear in $\theta$, $-a\sqrt{\theta^2/\alpha+\lambda^2r^2}$ is linear in $a$ and concave in $\theta$, and $\frac\lambda2a^2$ and $-\frac{\theta^2}2$ are strictly convex and strictly concave. The stationarity equations are
\begin{gather*}
\partial_\theta\cG=\sqrt{a^2+\sigma^2}-\frac{a\theta/\alpha}{\sqrt{\theta^2/\alpha+\lambda^2r^2}}-\theta=0,\\
\partial_a\cG=\frac{\theta a}{\sqrt{a^2+\sigma^2}}-\sqrt{\theta^2/\alpha+\lambda^2r^2}+\lambda a=0.
\end{gather*}
Let $(a,\theta)\in(0,\infty)^2$ solve them and define $\chi$ by $1+\chi=\sqrt{a^2+\sigma^2}/\theta$. The first equation says $\theta\chi=(a\theta/\alpha)/\sqrt{\theta^2/\alpha+\lambda^2r^2}$, so $\chi>0$ and $\sqrt{\theta^2/\alpha+\lambda^2r^2}=a/(\alpha\chi)$. Dividing the second equation by $a$ and inserting this gives $\frac1{1+\chi}-\frac1{\alpha\chi}+\lambda=0$, which is \eqref{eq:ridge:chi}, so $\chi=\chi_\alpha(\lambda)$. Squaring $\sqrt{\theta^2/\alpha+\lambda^2r^2}=a/(\alpha\chi)$, inserting $\theta^2=(a^2+\sigma^2)/(1+\chi)^2$ and multiplying by $\alpha^2\chi^2$ gives $\nu(a^2+\sigma^2)+r^2(\alpha\lambda\chi)^2=a^2$, that is, $a^2=\ridgeR(\alpha,\lambda)$, and then $\theta^2=\cT(\alpha,\lambda)$. So there is at most one stationary point. Conversely, with $\chi=\chi_\alpha(\lambda)$, $a_*=\ridgeR^{1/2}>0$ and $\theta_*=\sqrt{a_*^2+\sigma^2}/(1+\chi)$, the computation runs backward and both equations hold; a stationary point of a function convex in $a$ and concave in $\theta$ is a saddle point. This proves part~1.

For part~2, \cref{comp:ridge:saddle} shows that \eqref{eq:ridge:saddle} has exactly one solution with $\chi>0$ (the first equation determines $\chi$, the second $m=r^2(1-\alpha\lambda\chi)$, and the third is linear in $q$ with coefficient $1-\nu\ne0$) and that $q-2m+r^2=\ridgeR$. In the auxiliary problem, $\sqrt{\theta_*^2/\alpha+\lambda^2r^2}=a_*/(\alpha\chi)$ turns the overlap of step~5 into $r^2-r^2\alpha\lambda\chi=m$, and the squared component along $h$ into $\frac{\theta_*^2}\alpha\alpha^2\chi^2=\nu(a_*^2+\sigma^2)=\ridgeR-r^2(\alpha\lambda\chi)^2=q-m^2/r^2$.
\end{proof}

\textcite[Remark~5.2.2]{ThrampoulidisAbbasiHassibi2018} solve the same two equations for ridge regression with a Gaussian design, with their $\kappa$ for our $\chi$.\sidenote{Their Remark~5.2.3 notes that the optimally tuned ridge attains the minimum mean squared error when the signal is Gaussian: the Bayes reading of $\lambda^*$ in \cref{sec:ridge:physics}.} By \cref{ex:ridge:energy}, $-\Psi$ at the replica saddle point is also $\alpha\,\cG(a_*,\theta_*)$.

\begin{proofsketch}[Proof sketch of \cref{thm:ridge:risk}(i) by the CGMT]
The plan is to localize the minimizer of the auxiliary problem and transfer the localization. Since $\norm{(X^\top X+n\lambda I)^{-1}X^\top}\le(4n\lambda)^{-1/2}$ and $\norm y^2/n\to r^2+\sigma^2$, the minimizer and the maximizing $v$ lie in fixed balls with probability tending to one, so the CGMT applies on compact convex sets. Conditionally on $\varepsilon$, the auxiliary objective converges to $\cG$ uniformly on compacts; since $\cG$ is strictly convex in $a$, restricting it to $\abs{\norm w-a_*}\ge\epsilon$ raises its value by a deterministic gap, and the localization part of the CGMT \parencite[Theorem~6.1(iii)]{ThrampoulidisAbbasiHassibi2018} transfers the gap, so that $\norm{\hat\beta_\lambda-\beta^*}\to a_*$. Omitted: the uniform convergence and the verification of the localization hypotheses, which \textcite[Section~4 and Theorem~4.1]{ThrampoulidisAbbasiHassibi2018} carry out for general convex losses and regularizers.
\end{proofsketch}

\paragraph{What the general theorems add.} \textcite[Theorem~1]{HastieEtAl2022} prove part~(ii) for designs with independent entries with a finite moment of order $4+\eta$, and give non-asymptotic deterministic equivalents for general covariance (Theorems~2--3 and~5--6); \textcite{ChengMontanari2024} give dimension-free ones, and the formula is the kernel learning curve that \textcite{CanatarEtAl2021} derived by replicas (\cref{ch:kernels}). \textcite[Theorem~2.1]{DobribanWager2018} prove that for a random isotropic teacher the optimal ridge is $\sigma^2/(\alpha r^2)$ for every covariance, the rigorous form of Krogh and Hertz's remark. \textcite[Theorems~4 and~6]{BartlettEtAl2020} characterize, up to constants, when the interpolator's risk tends to zero; the isotropic proportional regime fails their conditions, as part~(ii) requires. And \textcite[Theorem~7]{HastieEtAl2022} prove that cross-validation and \term{generalized cross-validation}\index{generalized cross-validation} \parencite{GolubHeathWahba1979} select the optimal ridge, which in this chapter's language is the replica formula for the training error (\cref{ex:ridge:gcv}).

\section{Measuring it}\label{sec:ridge:measuring}

\begin{figure}[t]
\begin{fullwidth}
\centering
\begin{tikzpicture}
\begin{groupplot}[group style={group size=2 by 1,horizontal sep=17mm},rplot,width=82mm,height=52mm,
  xmode=log,xmin=0.23,xmax=5.4,xtick={0.25,0.5,1,2,4},xticklabels={$\frac14$,$\frac12$,$1$,$2$,$4$},
  xlabel={load $\alpha=n/d$},unbounded coords=jump]
\nextgroupplot[ylabel={excess risk $R$},ymin=0,ymax=2.2,ytick={0,0.5,1,1.5,2},title={(a) excess risk},
  legend to name=ridgecurveslegend,legend columns=3,legend style={/tikz/every even column/.append style={column sep=3mm}}]
\addplot[guide,forget plot] coordinates {(1,0) (1,2.2)};
\addplot[physpred] table[x=alpha,y=R0] {figures/data/09-ridge-curves.dat};
\addplot[proofpred] table[x=alpha,y=R0] {figures/data/09-ridge-curves.dat};
\addplot[physpred,forget plot] table[x=alpha,y=R1] {figures/data/09-ridge-curves.dat};
\addplot[proofpred,forget plot] table[x=alpha,y=R1] {figures/data/09-ridge-curves.dat};
\addplot[physpred,forget plot] table[x=alpha,y=Rs] {figures/data/09-ridge-curves.dat};
\addplot[proofpred,forget plot] table[x=alpha,y=Rs] {figures/data/09-ridge-curves.dat};
\addplot[measured,ebars] table[x=alpha,y=R0,y error plus=R0h,y error minus=R0h] {figures/data/09-ridge-sim.dat};
\addplot[measured2,ebars] table[x=alpha,y=R1,y error plus=R1h,y error minus=R1h] {figures/data/09-ridge-sim.dat};
\addplot[measured3,ebars] table[x=alpha,y=Rs,y error plus=Rsh,y error minus=Rsh] {figures/data/09-ridge-sim.dat};
\legend{replica prediction,{theorem (random matrices)},simulation $\lambda\rightarrow0$,simulation $\lambda=0.01$,simulation $\lambda=\lambda^*$}
\node[font=\scriptsize,text=physics,anchor=east] at (axis cs:0.78,1.75) {$\lambda\to0$};
\node[font=\scriptsize,text=physics,anchor=north east] at (axis cs:0.47,0.6) {$\lambda=\lambda^*$};
\nextgroupplot[ylabel={training error $T$},ymin=0,ymax=0.24,ytick={0,0.05,0.1,0.15,0.2},title={(b) training error},
  yticklabel style={/pgf/number format/.cd,fixed,precision=2}]
\addplot[guide,forget plot] coordinates {(1,0) (1,0.24)};
\addplot[physpred,forget plot] table[x=alpha,y=T0] {figures/data/09-ridge-curves.dat};
\addplot[proofpred,forget plot] table[x=alpha,y=T0] {figures/data/09-ridge-curves.dat};
\addplot[physpred,forget plot] table[x=alpha,y=T1] {figures/data/09-ridge-curves.dat};
\addplot[proofpred,forget plot] table[x=alpha,y=T1] {figures/data/09-ridge-curves.dat};
\addplot[physpred,forget plot] table[x=alpha,y=Ts] {figures/data/09-ridge-curves.dat};
\addplot[proofpred,forget plot] table[x=alpha,y=Ts] {figures/data/09-ridge-curves.dat};
\addplot[measured,ebars,forget plot] table[x=alpha,y=T0,y error plus=T0h,y error minus=T0h] {figures/data/09-ridge-sim.dat};
\addplot[measured2,ebars,forget plot] table[x=alpha,y=T1,y error plus=T1h,y error minus=T1h] {figures/data/09-ridge-sim.dat};
\addplot[measured3,ebars,forget plot] table[x=alpha,y=Ts,y error plus=Tsh,y error minus=Tsh] {figures/data/09-ridge-sim.dat};
\node[font=\scriptsize,text=physics,anchor=south] at (axis cs:1.5,0.15) {$\lambda=\lambda^*$};
\node[font=\scriptsize,text=physics,anchor=north] at (axis cs:2.6,0.12) {$\lambda\to0$, $\lambda=0.01$};
\end{groupplot}
\end{tikzpicture}

\vspace{1mm}
\pgfplotslegendfromname{ridgecurveslegend}
\end{fullwidth}
\caption{Double descent and its cure, with $r^2=1$ and $\sigma^2=0.25$. Curves: the replica prediction (carnelian) and \cref{thm:ridge:risk} (lapis dashed), which coincide, for the minimum-norm interpolator ($\lambda\to0$, divergent at $\alpha=1$), a small ridge $\lambda=0.01$ (the middle curve in (a)) and the optimal ridge $\lambda^*=\sigma^2/(\alpha r^2)$. Points: means over $8$ independent data sets with $n=1000$ and $d=\mathrm{round}(n/\alpha)$ at $16$ loads, with Student-$t$ $95\%$ intervals (seeds $90000+100k+j$ for load $k$ and data set $j$). All $48$ intervals in (a) contain the curves; in (b), $45$ of $48$ do (below $\alpha=1$ the interpolator's training error is exactly $0$). Scripts: \texttt{figures/scripts/09-ridge-routes.py} and \texttt{09-ridge-sim.py}.}\label{fig:ridge:curves}
\end{figure}

Three scripts measure the chapter's claims. \texttt{09-ridge-routes.py} computes $\ridgeR$ at $21$ pairs $(\alpha,\lambda)$ from the replica equations \eqref{eq:ridge:saddle} by fixed-point iteration, from \eqref{eq:ridge:risk}, from the Stieltjes form of \cref{rem:ridge:mp}, and from the auxiliary problem \eqref{eq:ridge:ao} by nested optimization. The first three agree to $6\times10^{-14}$, the last to $3\times10^{-8}$, its optimizer's tolerance; by \cref{prop:ridge:identity} this checks the programs, not the physics.

\texttt{09-ridge-sim.py} draws data sets with $n=1000$ at $16$ loads, with the teacher uniform on the unit sphere, noise variance $0.25$ and $8$ data sets per load, in about $85$ seconds.\sidenote{The estimators are computed by Cholesky solves of the smaller Gram matrix; for $d>n$ the push-through identity gives $\hat\beta_\lambda=X^\top(XX^\top+n\lambda I)^{-1}y$.} All $48$ risk intervals in \cref{fig:ridge:curves} contain the curves, and $29$ of the $32$ training-error intervals at $\lambda=0.01$ and $\lambda^*$ contain $\cT$ (the misses, at $\alpha=0.6$ and $0.7$, are $6$--$8\%$ low). The susceptibility, measured as the resolvent trace $\chi_n=\tr(X^\top X+n\lambda^*I)^{-1}$ of \cref{prop:ridge:chi}, is the best-measured quantity: over the $16$ loads its relative half-widths are $4\times10^{-5}$ to $4\times10^{-3}$, and its mean differs from $\chi$ by at most $0.3\%$, mostly upward, a bias of order $1/n$ that the narrow intervals detect at $4$ loads. The peak is the hardest place to measure: at $\alpha=0.9$ and $1.1$ the interpolator's risk is $2.37\pm0.30$ and $2.51\pm0.31$, against limits $2.35$ and $2.50$ and exact finite-$n$ means (\cref{lem:ridge:wishart}) $2.37$ and $2.53$, relative half-widths of $13\%$ and $12\%$ against $2\%$ and $6\%$ for the optimal ridge, because $\tr(X^\top X)^{-1}$ has a heavy tail near $n=d$ (\cref{ex:ridge:tails}).

\texttt{09-ridge-universal.py} repeats the simulation at $n=500$ with independent $\pm1$ entries and with Gaussian ones ($12$ data sets, $\alpha\in\{0.5,2\}$, $\lambda\in\{0,\lambda^*\}$). Seven of the eight risk intervals contain the limits (the miss is Gaussian, $1.4$ half-widths high), and the conditional means $\E[R\mid X]$, computable for every design because the teacher is uniform, agree with the limits within $0.7\%$ for both designs ($0.6406\pm0.0001$ Gaussian and $0.6403\pm0.0001$ for $\pm1$ entries at $\alpha=0.5$ and $\lambda^*$, limit $0.6404$), although every Gaussian step of our proof fails for $\pm1$ entries (\cref{ex:ridge:universal}).

\section{The status board}\label{sec:ridge:status}

As of October 2026 every prediction in the physics box is a theorem for isotropic Gaussian designs, and its ridgeless part also for designs with independent entries; the replica computation itself remains a heuristic.

\begin{statusboard}
\statusrow{Replica-symmetric equations and formulas for $R$, $T$}{\heuristicstatus}{\cref{sec:ridge:physics}; finite temperature: \textcite{DunmurWallace1993}}
\statusrow{Ridge risk and training error for $\lambda>0$, isotropic Gaussian design}{\proved}{\cref{thm:ridge:risk}(i); \textcite[Cor.~5]{HastieEtAl2022}; \textcite[Thm.~2.1]{DobribanWager2018}}
\statusrow{Ridgeless limits, divergence at $\alpha=1$, exchange of $\lambda\to0$ and $n\to\infty$}{\proved}{\cref{thm:ridge:risk}(ii); entries with $4+\eta$ moments: \textcite[Thm.~1]{HastieEtAl2022}}
\statusrow{Optimal ridge $\lambda^*=\sigma^2/(\alpha r^2)$, risk $\sigma^2\chi$, no peak}{\proved}{\cref{thm:ridge:risk}(iii); every $\Sigma$, random teacher: \textcite[Thm.~2.1]{DobribanWager2018}; finite $n$: \textcite[Prop.~1]{Dicker2016}}
\statusrow{Replica equations $=$ auxiliary-problem equations $=$ resolvent equation}{\proved}{\cref{prop:ridge:chi,prop:ridge:identity}; same system: \textcite[Rem.~5.2.2]{ThrampoulidisAbbasiHassibi2018}}
\statusrow{Cross-validation and GCV select the optimal ridge}{\proved}{isotropic, \textcite[Thm.~7]{HastieEtAl2022}}
\statusrow{The learning curve for general covariance (replicas: \textcite{CanatarEtAl2021})}{\provedunder}{bounded spectrum, moments \parencite[Thms.~2--3, 5--6]{HastieEtAl2022}; concentration \parencite{ChengMontanari2024}}
\statusrow{Conditions for benign overfitting of the interpolator}{\proved}{up to constants, \textcite[Thms.~4, 6]{BartlettEtAl2020}}
\statusrow{The deep-network peak is a jamming transition}{\conjectured}{numerics, \textcite{SpiglerEtAl2019}; \textcite{GeigerEtAl2019}}
\end{statusboard}

\section{Where the translation fails}\label{sec:ridge:fails}

\paragraph{What was assumed.} Both the replica computation and the theorem take isotropic Gaussian inputs, a well-specified linear model and noise independent of the inputs, which for the proof has a finite fourth moment, used only for the variance term. The proof uses Gaussianity in rotation invariance, the Poincaré inequality and the inverse Wishart mean; the comparison route needs it throughout. The exchange of limits is not assumed: part~(ii) proves it.

\paragraph{Where the physics is wrong, or not yet right.} In the proportional limit nothing in the physics box is wrong; what fails is its reading at finite size and for a single teacher. Self-averaging fails near the threshold: for Gaussian designs and $\sigma>0$ the interpolator's risk has infinite mean when $\abs{n-d}\le1$ and infinite variance when $\abs{n-d}\le3$ (\cref{ex:ridge:tails}). Krogh and Hertz's spectrum-independent optimal decay holds on average over an isotropic teacher, as their derivation and \textcite{DobribanWager2018} assume; for a fixed teacher and an anisotropic covariance the optimal ridge depends on how the teacher aligns with the covariance, and it can be zero or negative \parencite{KobakEtAl2020}. And the proposal that the interpolation peak of deep networks is a jamming transition \parencite{SpiglerEtAl2019,GeigerEtAl2019} is a conjecture supported by numerics, which the linear model, whose peak comes from small eigenvalues of $X^\top X/n$, neither confirms nor refutes.

\paragraph{Open problems.} The isotropic model is closed; its neighbors are not. First, real data: \textcite{LoureiroEtAl2022} conjecture that the Gaussian-design formula, with the feature covariance, predicts the ridge risk of generic feature maps of realistic data (their Conjecture~1; its analogue for the classification error fails on real images), and \textcite{ChengMontanari2024} prove it under concentration hypotheses that real data need not satisfy (\cref{ch:kernels}). Second, for random-features ridge regression the exchange of limits proved in part~(ii) is a conjecture \parencite[Remark~1]{MeiMontanari2022}. Third, a question of our own: here the integer moments $\E Z^k$ are explicit random-matrix integrals, so this is a natural model on which to test whether their continuation to $k\to0$ reproduces the free entropy (\cref{ex:ridge:replicatrick}); we know of no such computation.

\begin{exercises}
\item \consolidation (The optimal ridge on the Nishimori line.) (a) Prove the Bayes reading of $\hat\beta_{\lambda^*}$ in \cref{sec:ridge:physics}, and deduce $m=q$ from the Nishimori identity. (b) Give a second proof that $\alpha\mapsto\ridgeR(\alpha,\lambda^*(\alpha))$ is non-increasing, without computing it. \hint{for fixed $d$ the Bayes risk does not increase with $n$.}\label{ex:ridge:nishimori}

\item \consolidation (Facts in the load.) Let $\mathrm{SNR}=r^2/\sigma^2$. Show that least squares ($\alpha>1$) beats the null estimator $\hat\beta=0$, of risk $r^2$, exactly when $\alpha>1+1/\mathrm{SNR}$, and the interpolator ($\alpha<1$) exactly when $\alpha<1-1/\mathrm{SNR}$; and that if $\sigma<r$ the interpolator's risk on $(0,1)$ is smallest at $\alpha=1-\sigma/r$, where it equals $\sigma(2r-\sigma)$.\label{ex:ridge:facts}

\item \consolidation (Finite-size optimality.) For $\beta^*=ru$ with $u$ uniform and independent of $X$, prove the formula for $\E[R(\hat\beta_\lambda)\mid X]$ stated after the proof of \cref{thm:ridge:risk} and its minimization at $\lambda=\sigma^2d/(nr^2)$. Why does only the second moment of the prior matter?\label{ex:ridge:dicker}

\item \byhand (Energy and auxiliary value.) Show that at the replica saddle point $-\Psi=\frac{\alpha\lambda q}2+\frac\alpha2\cT$, the minimum of $H/d$, and that this equals $\alpha\,\cG(a_*,\theta_*)$. \hint{use the third equation of \eqref{eq:ridge:saddle} to eliminate $(q-m^2/r^2)/\chi$.}\label{ex:ridge:energy}

\item \byhand (Gordon for least squares.) For $\lambda=0$ and $\alpha>1$ solve $\min_a\max_\theta\{\theta\sqrt{a^2+\sigma^2}-a\theta/\sqrt\alpha-\theta^2/2\}$ by hand and recover $R=\sigma^2/(\alpha-1)$ and $T=\sigma^2(1-1/\alpha)$. What goes wrong for $\alpha<1$?\label{ex:ridge:ols}

\item \byhand\hard (Heavy tails at the threshold.) Let $\sigma>0$ and the noise be Gaussian. Show that $\E R(\hat\beta_0)<\infty$ exactly when $\abs{n-d}\ge2$, and $\Var R(\hat\beta_0)<\infty$ exactly when $\abs{n-d}\ge4$. \hint{$\E[(\chi^2_m)^{-2}]=1/((m-2)(m-4))$ for $m>4$ and is infinite for $m\le4$; use $(\tr W^{-1})^2\ge(W^{-1})_{11}^2$ and, for the converse, $\tr W^{-2}\le(\tr W^{-1})^2\le d\sum_j(W^{-1})_{jj}^2$.}\label{ex:ridge:tails}

\item \translate (The susceptibility of a Gaussian Gibbs measure.) Show that the Gibbs measure of \eqref{eq:ridge:energy} at inverse temperature $b$ is $\Normal(\sqrt d\,\hat\beta_\lambda,(bH'')^{-1})$ with $H''=X^\top X/d+\alpha\lambda I$, so that $b(\gibbs{Q}-\gibbs{q})=\tr(X^\top X+n\lambda I)^{-1}$ exactly. Which replica steps does this make unnecessary?\label{ex:ridge:gibbs}

\item \translate (The training error is generalized cross-validation.) Let $H_\lambda=X(X^\top X+n\lambda I)^{-1}X^\top$. Show that $\tr H_\lambda/n=\frac dn-\lambda\chi_n\to\chi/(1+\chi)$, and deduce that $\mathrm{GCV}(\lambda)=T(\hat\beta_\lambda)/(1-\tr H_\lambda/n)^2\to\ridgeR(\alpha,\lambda)+\sigma^2$. Which replica formula is this?\label{ex:ridge:gcv}

\item \measurement (Universality.) Run \texttt{09-ridge-universal.py} and compare both designs with the limits. Why does the conditional mean of \cref{ex:ridge:dicker}, valid for every design, fluctuate far less than the risk? Which steps of the proof of \cref{thm:ridge:risk} fail for $\pm1$ entries?\label{ex:ridge:universal}

\item \research (The replica trick on a solvable model.) For the energy \eqref{eq:ridge:energy}, $Z=e^{-b\min H}(2\pi/b)^{d/2}\det(H'')^{-1/2}$. Compute $\lim_d\frac1{dk}\log\E Z^k$ for integer $k$, a large-deviation problem for $(\min H/d,\frac1d\log\det H'')$, and decide whether its continuation to $k\to0$ gives $\lim_d\frac1d\E\log Z$.\label{ex:ridge:replicatrick}
\end{exercises}

\notesandrefs
Ridge regression is due to \textcite{HoerlKennard1970} and generalized cross-validation to \textcite{GolubHeathWahba1979}. The prehistory of double descent, from the pseudo-inverse rule of 1989--1990, is told by \textcite{LoogEtAl2020}, its name is from \textcite{BelkinEtAl2019}, and \textcite{MeiMontanari2022} proved it for random-features regression (\cref{ch:kernels}).

The rigorous theory of the isotropic case was built in steps: \textcite{MarchenkoPastur1967} for the spectrum, \textcite{Dicker2016} for the optimally tuned risk, \textcite{DobribanWager2018} for general covariance with a random teacher, and \textcite{HastieEtAl2022} for the ridgeless limits, general covariance and cross-validation. Gordon's comparison \parencite{Gordon1988} gives the second proof through the CGMT. The leave-one-out method, the rigorous form of the cavity method, was used by \textcite{ElKaroui2013} for ridge-regularized robust regression.

The replica derivation of \cref{sec:ridge:physics} is the standard zero-temperature computation, and the proof of \cref{thm:ridge:risk} is assembled from standard tools: the leave-one-out proof of \cref{prop:ridge:chi} is a version of the usual Schur-complement route to the Marchenko--Pastur law, and the treatment of the ridgeless limit by a monotone comparison and the exact inverse Wishart mean is, as far as we know, the book's own. \Cref{prop:ridge:identity} and the dictionary of the stone are the book's formulation; for the Gaussian covariate model, which contains ridge regression, \textcite[Appendix~C]{LoureiroEtAl2022} showed that the replica and Gordon equations coincide.

\paragraph{Translating the literature.} \textcite{HastieEtAl2022} write $p$ for our $d$ and $\gamma=p/n=1/\alpha$; their objective $\frac1n\norm{y-Xb}^2+\lambda\norm b^2$ has our minimizer, their $m(z)$ is the Marchenko--Pastur Stieltjes transform with $m(-\lambda)=\alpha\chi$, and in their Theorem~7 the letter $\alpha$ means $r^2/(\sigma^2\gamma)$. \textcite{DobribanWager2018} use $\gamma=p/n$, our $\lambda$, noise variance $1$, signal strength $\alpha^2$ for $r^2/\sigma^2$, a risk that includes the noise, and a companion transform $v$ with $\lambda v(-\lambda)=1/(1+\chi)$. \textcite{Dicker2016} writes $(X^\top X+\frac d{t^2}I)^{-1}X^\top y$, so that $\lambda=d/(nt^2)$. \textcite{ThrampoulidisAbbasiHassibi2018} have $m$ measurements of $n$ unknowns with $A_{ij}\sim\Normal(0,1/n)$, so their $\delta$ is our $\alpha$ and their $x$ our $w=\sqrt d\,\beta$; their ridge parameter is $\alpha\lambda$, their $\kappa$ our $\chi$, their $\alpha$ our $\sqrt R$, and their $\sigma_x^2,\sigma_z^2$ our $r^2,\sigma^2$. \textcite{KroghHertz1991} write $N$ inputs, $p=\alpha N$ patterns, the decay $\alpha\lambda$, $\sigma_u^2$ (their $u^2$) for $r^2$, $F$ for $R$, and $A$ for $X^\top X/d$. \textcite{CanatarEtAl2021} write $\kappa$ for $n\lambda_{\mathrm{eff}}=n\lambda(1+\chi)$.

\printchapterbib
\end{refsection}
